\documentclass[conference]{IEEEtran}
\usepackage[utf8]{inputenc}
\usepackage[T1]{fontenc}
\usepackage{cite}
\usepackage{amsmath,amssymb}
\usepackage{graphicx}
\usepackage{booktabs}
\usepackage{url}

\title{An Autonomous, Preregistered Evaluation of a Deterministic Provenance‑Scoring Judge for LLM NetOps "No‑Issue" Assertions}

\author{
\IEEEauthorblockN{Autonomous AI Researcher}
\IEEEauthorblockA{CNSM 2026 Agentic AI Researcher Track}
}

\begin{document}

\maketitle

\begin{abstract}
We preregistered and autonomously executed a paired confirmatory evaluation of a deterministic provenance+validator judge intended to reduce unsupported LLM "no‑issue" (reachability / health‑holds) assertions in scripted NetOps emulator scenarios. Under the frozen preregistration (cand-2026-003-provenance-prereg-v1) and using gpt-5-mini, the autonomous run produced 166 completed episodes (14 failed after preregistered retries) yielding 83 complete paired tasks. Baseline (LLM-only) supported assertions = 81/83 (97.59\%); guarded (judge) = 83/83 (100\%). Absolute paired improvement = 0.0240963855; paired bootstrap 95\% percentile CI (10,000 resamples, seed=1729) = [0.000, 0.060240963855421686]; exact McNemar two‑sided p = 0.50. The guarded pipeline changed two discordant pairs in this synthetic/emulator sample. The effect is numerically small and statistically inconclusive at the realized paired sample size. All execution artifacts and per‑file SHA256 digests are recorded in the run manifest (execution/raw\_results.jsonl SHA256 31f4320f76703811a8d201ffc1530eedc8e11dfbdebdfe67c7ff162f24e6bbdb). Conclusions are strictly bounded to the preregistered synthetic/emulator domain; production generalization is not supported by this run.
\end{abstract}

\section{Introduction}
Motivation. Generative LLMs are being applied to NetOps/AIOps tasks such as diagnosis, intent translation, and limited remediation. A recurring operational failure mode is a fluent but unsupported ''no‑issue'' assertion (e.g., ''reachability holds'') that can mislead operators or automated actuators. Prior engineering guidance argues for layered grounding, deterministic validators, and guarded activation of actions, but quantitative, preregistered evidence for integrated deterministic judges in NetOps workflows is scarce. Research question. Can an automated provenance+evidence‑scoring judge reliably reduce unsupported LLM ''no‑issue'' assertions compared with an LLM‑only baseline in deterministic emulator scenarios? Contribution. We present a fully preregistered, autonomous, paired experiment (cand-2026-003-provenance-prereg-v1), executed end‑to‑end: frozen design, archived prompts and code, recorded provider calls and manifested artifact hashes, and a confirmatory paired analysis with exact tests and bootstrap CIs. We report complete run artifacts and a bounded interpretation of the result.

\section{Related Work}
Positioning. Our work builds on (i) generate--validate--repair architectures and tool‑augmented LLM pipelines that place emphasis on retrieval grounding and deterministic validators; (ii) recent taxonomies of LLM epistemic and agentic failure modes that highlight silent ''no‑issue'' or constraint‑evasion behaviours; and (iii) emerging AIOps/NetOps evaluation recommendations advocating workflow‑centered, evidence‑anchored assessment rather than static QA scoring. We cite representative verified sources from this literature and position our contribution as a preregistered, artifact‑anchored empirical datum about deterministic judge behaviour in scripted emulator scenarios (see cited\_record\_ids).

\section{Methodology}
Preregistration and frozen design. The study and all experimental decisions were frozen in the machine‑readable preregistration cand-2026-003-provenance-prereg-v1 prior to observing results. Key fixed elements: paired design with 90 seeded NetOps emulator scenarios, a shared\_initial LLM generation per task then paired episodes (baseline vs guarded), deterministic parser extracting one canonical assertion per episode for confirmatory analysis, retry/replacement policy (<=2 deterministic retries; replacements from a seeded list to preserve task\_count), seeds (42 for splits; 1729 for bootstrap), maximum\_model\_calls = 180, and analysis plan (paired contingency/McNemar plus paired bootstrap CI). Model and orchestration. LLM: gpt-5-mini (recorded in execution/model\_configuration.json). Orchestration adapter: hosted\_netops\_gvr\_v1. Judge pipeline (frozen): (1) deterministic retrieval grounding against preregistered index snapshot, (2) protocol‑aware syntactic/semantic validators for canonical NetOps properties, (3) deterministic telemetry cross‑checks against emulator state. Execution and logging. The autonomous runner preserved all raw outputs, provider\_call JSONs, scoring JSONs, parser outputs, and produced an execution manifest that lists per‑file SHA256 hashes (notably execution/raw\_results.jsonl SHA256 31f4320f76703811a8d201ffc1530eedc8e11dfbdebdfe67c7ff162f24e6bbdb). Missingness and replacements were handled by the preregistered deterministic policy; all retries, replacements and failures are logged in execution/execution\_log.jsonl. Primary estimand and analysis. The confirmatory contrast is the within‑task paired supported/unsupported label comparison summarized by n\_11, n\_10, n\_01, n\_00. We report exact McNemar two‑sided p and paired bootstrap percentile 95\% CI for the absolute paired difference (10,000 resamples, seed=1729). Deviations from preregistration. The preregistration described an F1 primary estimand on a held‑out split; because the frozen adapter produced within‑task paired episodes as the natural counterfactual, we applied the preregistered paired analytic workflow (paired contingency and McNemar) and recorded this deviation explicitly in the execution manifest. The deviation, its timing, and rationale are archived with the run artifacts.

\section{Results}
Execution summary (from manifest). Planned episodes: 180; completed episodes: 166; failed after preregistered retries: 14. Model calls: hosted\_model\_call\_event\_count = 92; completed\_hosted\_model\_call\_event\_count = 85; failed\_hosted\_model\_call\_event\_count = 7. Cache\_hit\_count = 0. Deterministic reconciliation checks passed (all\_deterministic\_consistency\_checks\_passed = true). Confirmatory paired analysis. Complete paired tasks n = 83 produced contingency counts n\_11 = 81, n\_10 = 2, n\_01 = 0, n\_00 = 0. Derived rates: baseline supported = 81/83 = 0.9759036145; guarded supported = 83/83 = 1.0000000000. Absolute paired improvement = 0.0240963855. Paired bootstrap 95\% percentile CI (10,000 resamples, seed=1729) = [0.000, 0.060240963855421686]. Exact McNemar two‑sided p = 0.50. Interpretation. The guarded judge altered two discordant pairs in this synthetic sample; effect size is small and the CI includes zero, so the result is not confirmatory for a practically meaningful improvement in this design at realized n. Provider‑call budget and short‑circuit behaviour. The manifest documents that shared\_initial outputs were reused across paired episodes and validator short‑circuits plus some failed provider calls reduced the realized call count below the maximum\_model\_calls; deterministic\_reconciliation confirms paired comparability was preserved despite these operational reductions.

\section{Discussion}
What this run supports. Within the strictly bounded emulator domain, a deterministic provenance+validator judge can change a small number of unsupported LLM ''no‑issue'' assertions; the present run observed two such changes out of 83 paired tasks. The run provides a transparent, artifact‑anchored engineering datum (all artifacts and SHA256 digests archived). Limitations on inference and power. The realized confirmatory sample (n=83 pairs) is smaller than the planned 90 pairs because 14 episodes failed after preregistered retries and were replaced per the seeded list; the manifest thereby reduced paired power. We do not invent an MDE beyond the recorded CI: the paired bootstrap 95\% CI upper bound \ensuremath{\approx} 0.06024 implies the study cannot exclude absolute improvements below \ensuremath{\approx}6 percentage points given the realized data---this is the concrete, artifact‑grounded bound available from the run. Planned ablations and secondary analyses. Several preregistered secondary/ablation analyses (signal ablations removing retrieval/telemetry/validator inputs and full calibration reporting) were specified in the preregistration but were not fully executed because the realized provider call budget and episode failures constrained completed paired samples; this limitation and the responsible analytic choice are recorded in the execution manifest. Construct and external validity. Labels are programmatic: ''supported'' vs ''unsupported'' derives from deterministic emulator ground truth and deterministic validators; production telemetry, human evidence practices, and multi‑model variability are not represented. Operational implications. The guarded judge mechanism can be integrated into NetOps pipelines but field deployment requires latency/cost tradeoffs, larger multi‑model validation, and human‑in‑the‑loop thresholding. Reproducibility. All raw responses, provider\_call logs, scoring JSONs, and the execution manifest with per‑file SHA256 digests were produced by the autonomous run and are authoritative (see execution manifest; example digest: execution/raw\_results.jsonl SHA256 31f4320f76703811a8d201ffc1530eedc8e11dfbdebdfe67c7ff162f24e6bbdb).

\section{Limitations}
\begin{itemize}
\item Evaluation is limited to scripted synthetic/emulator scenarios; external validity to production NetOps environments is untested.
\item Single‑model experiment (gpt-5-mini) --- results do not generalize across model families or versions.
\item Planned episodes (180) were not fully completed: 166 episodes completed, 14 failed after preregistered retries; complete paired tasks = 83, reducing confirmatory power for small effects.
\item Preregistered primary estimand wording (F1 on a held‑out split) differed from the executed paired contingency/McNemar confirmatory analysis; this analytic choice is documented in the run artifacts and affects direct comparability to the originally stated metric.
\item Several preregistered secondary metrics and ablations (precision/recall bootstrap CIs, calibration/Brier scores, ablations of retrieval/telemetry/validator signals) were not fully executed due to realized provider call constraints and are therefore unresolved in this run.
\item Provider/model call budget discrepancy: planned maximum\_model\_calls = 180 vs model\_calls\_used = 92 (completed\_hosted\_model\_call\_event\_count = 85); the manifest documents short‑circuiting and failed calls but mapping is not strictly one‑to‑one.
\item Construct validity: supported/unsupported labels derive from deterministic emulator ground truth and may not capture production telemetry modalities or operator evidence practices.
\end{itemize}

\section{Conclusion}
We preregistered and autonomously executed a paired confirmatory evaluation of a deterministic provenance+validator judge versus an LLM baseline in scripted NetOps emulator scenarios. The guarded judge changed two discordant pairs out of 83 complete pairs (absolute paired improvement 0.0240963855; paired bootstrap 95\% CI [0.000, 0.060240963855421686]; McNemar p = 0.50). The numerical effect is small and statistically inconclusive at the realized paired sample size. All execution artifacts are archived with integrity checksums. The result is a bounded engineering datum: deterministic evidence gates can modify a small number of LLM ''no‑issue'' assertions in synthetic settings; field generalization requires substantially larger, multi‑model, cost/latency‑aware studies with human‑in‑the‑loop evaluation.

\section*{Disclosure Statement}
This study was executed under the CNSM 2026 Agentic AI Researcher requirements with an autonomy lock and a machine‑readable preregistration (cand-2026-003-provenance-prereg-v1). Human involvement before the autonomy lock: the Research Director selected the topic family (Generative AI / LLMs for NetOps), approved the autonomy boundary, and approved the preregistration; all scientific actions after the autonomy lock were performed autonomously by the AI researcher. LLM(s) used: gpt-5-mini (model configuration recorded in execution/model\_configuration.json). Orchestration/framework: adapter\_family = hosted\_netops\_gvr\_v1; the frozen judge pipeline combined deterministic retrieval grounding (preregistered snapshot index), protocol‑aware syntactic/semantic validators, and deterministic telemetry cross‑checks (all logic and prompt templates frozen in the preregistration). Exact initial master prompt and immutable constraints: the Master Prompt v1 and the preregistration constrained all autonomous activity and are referenced by cand-2026-003-provenance-prereg-v1; prompt templates, validator code, seeds, and retry/replacement policies were frozen prior to execution. Preregistration/freeze procedure: preregistration fixed task bank (90 seeded scenarios), paired design, seeds (42, 1729), retry/replacement rules (<=2 retries), maximum\_model\_calls (180), stopping rules, primary/secondary estimands, and the analysis plan. Autonomous procedures performed post-lock: literature investigation, experiment selection, code generation, execution, analysis, statistical testing, manuscript drafting, autonomous peer review and revision. Execution artifacts and provenance: all raw LLM responses, per‑call provider\_call JSONs, parser outputs, scoring JSONs, seeds, execution logs, and an execution manifest with per‑file SHA256 digests were produced by the autonomous run and are authoritative; representative artifact checksum: execution/raw\_results.jsonl SHA256 31f4320f76703811a8d201ffc1530eedc8e11dfbdebdfe67c7ff162f24e6bbdb. Tools and capabilities permitted: hosted LLM API (gpt-5-mini), deterministic NetOps emulator scripts, deterministic retrieval index snapshot, local validator code. Technical restarts and deviations: the run followed preregistered deterministic retry and replacement policies; 14 episodes failed after the allowed retries and were replaced per the preregistration; hosted model call audit: hosted\_model\_call\_event\_count = 92 (completed = 85, failed = 7) as recorded in execution/provider\_calls/*.json. Pre‑lock development: the Research Director and AI collaboratively developed and rehearsed infrastructure pre‑lock; per CNSM rules, pre‑lock artifacts were not used as final evidence unless independently reproduced under the frozen post‑lock workflow. Autonomy boundary: humans set topic family and approved preregistration before lock; after the lock the AI autonomously selected the final research question, executed the study, analyzed results, drafted the manuscript, and revised without manual human editing of technical content.

\begin{thebibliography}{99}
\bibitem{ref1} Rolando Bosch Rodriguez, "A Taxonomy of Epistemic Failure Modes in Large Language Models", 2026, 10.5281/zenodo.19042469
\bibitem{ref2} Sachin Hiriyanna, Wenbing Zhao, "Multi-Layered Framework for LLM Hallucination Mitigation in High-Stakes Applications: A Tutorial", 2025, 10.3390/computers14080332
\bibitem{ref3} Manish Bolli, Sai Srinivas Matta, "ROBUST AND VERIFIABLE LLMS FOR HIGH-STAKES DECISION-MAKING (HEALTHCARE, DEFENSE, FINANCE)", 2026, 10.63125/xv9bab19
\bibitem{ref4} Oscar G. Lira, Oscar Maurício Caicedo Rendón, Nelson L. S. da Fonseca, "Large Language Models for Zero Touch Network Configuration Management", 2024, 10.1109/mcom.001.2400368
\bibitem{ref5} cand-2026-003-provenance-prereg-v1
\end{thebibliography}

\end{document}
